\documentclass[11pt]{article}
\usepackage{explorer}
\newcommand{\ExplorerLanguage}{FR}
\newcommand{\ExplorerNumber}{09}
\newcommand{\ExplorerKind}{Solutions détaillées}
\begin{document}
\ExplorerTitle{Mathématiques expertes · géométrie}{Rotations et nombres complexes}
\ExplorerSection{Exercice 1}
Calculer \(i(2-3i)\), puis son module.\par
\(i(2-3i)=3+2i\) ; son module est \(\sqrt{13}\), comme celui de \(2-3i\).\par
\ExplorerSection{Exercice 2}
Calculer les images successives de \(1+i\) par quatre multiplications par i.\par
−\(1+i, -1-i, 1-i\), puis \(1+i\). Quatre quarts de tour redonnent le point initial.\par
\ExplorerSection{Exercice 3}
La transformation \(z'=i z+2 a\) quel centre ?\par
Le centre vérifie \(a=i a+2\), donc \(a=\frac{2}{1-i}=1+i\).\par
\ExplorerSection{Exercice 4}
Pour \(z'=a+i(z-a)\), prouver |\(z'-a|=|z-a|\).\par
\(z'-a=i(z-a)\), donc |\(z'-a|=|i||z-a|=|z-a|\).\par
\ExplorerSection{Problème : raisonnement complet}
On définit \(z_{0}=2\) et \(z_{n+1}=i z_{n}+2\). Trouver le point fixe a, puis une formule explicite de \(z_{n}\). Déterminer la période minimale de la suite et la somme \(z_{0}+z_{1}+z_{2}+z_{3}\). Interpréter géométriquement les quatre points.\par
\(a=1+i\) et \(z_{n+1}-a=i(z_{n}-a)\), donc \(z_{n}=1+i+i^{n}(1-i)\). Les valeurs sont \(2, 2+2i, 2i, 0\), puis 2. Elles sont distinctes sur les quatre premiers rangs : la période minimale est 4. Leur somme vaut \(4+4i=4a\). Les points forment un carré de centre (1,1), de côté 2 et de rayon \(\sqrt{2}\).\par
\end{document}
